\chapter{CONCLUSIONS AND FUTURE WORK}

\section{Summary of Contributions}

This thesis has presented CTU-SignBridge, a purpose-built platform that directly addresses the severe data-collection bottleneck hindering Vietnamese Sign Language (VSL) AI research. By transitioning from ad-hoc manual data collection to a structured, highly automated cloud architecture, the system provides a robust foundation for building the largest open-source VSL dataset to date. The key contributions of this research are:

\begin{enumerate}
    \item \textbf{Five-Layer Clean Architecture:} Implemented a strict separation of concerns within the backend, enabling highly maintainable code and independent testing of routers, schemas, services, repositories, and async workers.
    \item \textbf{MediaPipe-First Storage Strategy:} Successfully shifted the computational burden of landmark extraction to the client's browser via WebAssembly. This reduced per-clip storage requirements by 99.1\%, making large-scale community data collection economically viable on commodity hardware.
    \item \textbf{Star Schema Taxonomy Design:} Built a high-performance relational structure featuring three fixed depth levels and two satellite dimension tables, enabling complex faceted searches (by dialect, category, and gesture feature) with query latencies consistently under 5\,ms.
    \item \textbf{Secure Multi-Tenant IAM Architecture:} Integrated Keycloak for robust Single Sign-On (SSO) authentication and Casbin for granular Role-Based Access Control (RBAC). This ensures provably isolated Workspace and Project environments for independent research groups sharing the same infrastructure.
    \item \textbf{Resilient Asynchronous Synchronization:} Leveraged Celery and Redis to handle cloud storage operations (via MinIO), achieving a 100\% eventual success rate in data synchronization despite simulated network timeouts or worker failures.
\end{enumerate}

All predefined system Key Performance Indicators (KPIs) were met successfully: P95 API latency remained $\leq$ 100\,ms under high concurrent load, storage reduction exceeded 90\%, and the asynchronous task queue demonstrated $\geq$ 99\% reliability.

\section{Future Directions}

While the current platform establishes a solid data engineering pipeline, several avenues remain for future enhancement:

\begin{itemize}
    \item \textbf{Facial Expression and Full-Body Capture:} The database currently reserves the \texttt{requires\_face\_expression} flag. Future iterations should expand the MediaPipe integration to concurrently capture the 468 face-mesh landmarks, as facial micro-expressions are critical grammatical markers in VSL.
    \item \textbf{Active Learning Feedback Loop:} Integrate a real-time inference model that evaluates the quality of uploaded clips. If the model flags a clip as low-confidence or noisy, the system could automatically generate a task in the queue for the contributor to re-record the specific sign.
    \item \textbf{Federated Data Collection:} Extend the multi-tenancy model to support cross-workspace data federation. This would allow research teams to pool their independent datasets to train a global VSL model while still retaining strict ownership and access control over their original data.
    \item \textbf{Horizontal Auto-Scaling:} As user adoption grows, the current Docker Compose deployment should be migrated to a Kubernetes (K3s) cluster. Implementing Horizontal Pod Autoscaling (HPA) would ensure the API and Celery workers scale dynamically based on CPU and memory metrics during peak collection campaigns.
    \item \textbf{Progressive Web App (PWA) Offline Mode:} Enhance the React frontend to function as an offline-first PWA. Contributors in remote areas with unstable internet connectivity could record and temporarily cache landmark data in IndexedDB, which the system would then automatically synchronize once connectivity is restored.
\end{itemize}
